% Synthetisches Testdokument: Formelsammlung im Querformat mit eigenen Boxen,
% mehrspaltigem Satz, eigenen Makros und tabbing.
\documentclass[a4paper,9pt,landscape]{extarticle}
\usepackage[utf8]{inputenc}
\usepackage[T1]{fontenc}
\usepackage{amsmath,amssymb,bm}
\usepackage[margin=8mm]{geometry}
\usepackage{multicol}
\usepackage{tcolorbox}
\usepackage{tikz}

\newtcolorbox{formelbox}[1]{colback=gray!5, colframe=gray!60, title={#1}, fonttitle=\bfseries}
\newcommand{\teil}[1]{\textbf{#1)}~}
\DeclareMathOperator{\rot}{rot}
\setlength{\parindent}{0pt}

\begin{document}
\begin{multicols*}{3}
\section*{Elektromagnetische Felder}

\begin{formelbox}{Maxwell-Gleichungen}
$\rot \vec{E} = -\dfrac{\partial \vec{B}}{\partial t}$\quad $\operatorname{div} \vec{B} = 0$\\[2pt]
\teil{a} Induktion \teil{b} Quellenfreiheit
\end{formelbox}

\begin{formelbox}{Größen}
\begin{tabbing}
\hspace*{1.5cm} \= \kill
$\varepsilon_0$ \> Feldkonstante \\
$\mu_0$ \> Permeabilität
\end{tabbing}
\end{formelbox}

\begin{center}
\begin{tikzpicture}
\draw[->] (0,0) -- (2,0) node[right] {$x$};
\draw[->] (0,0) -- (0,1.5) node[above] {$y$};
\end{tikzpicture}
\end{center}

\columnbreak
\section*{Wellen}
Wellenzahl $k = 2\pi/\lambda$.
\end{multicols*}
\end{document}
